\chapter{User Guide}
\label{app:user-guide}

\section{Purpose}
\label{sec:user-guide-purpose}

\noindent
 This appendix explains how to run the \textbf{Smart Photos} prototype on a local machine. Also a deployed version can be accessed at \url{https://smart-photos.vercel.app}.
A concise, illustrated tour of every major flow is already given in Section~\ref{sec:translation}. 

\section{Prerequisites}

\begin{itemize}
  \item \textbf{Node.js} version \textbf{20.x} (LTS recommended). The project uses TypeScript tooling aligned with Node~20; older major versions may fail to install or run dependencies.
  \item \textbf{npm} (bundled with Node.js). The instructions below use \texttt{npm}
\end{itemize}

\section{Install and run}

\begin{enumerate}
  \item Open a terminal and change into the prototype folder:
  \begin{verbatim}
cd photo-revocation-prototype
  \end{verbatim}
  \item Install dependencies:
  \begin{verbatim}
npm install
  \end{verbatim}
  \item Start the development server:
  \begin{verbatim}
npm run dev
  \end{verbatim}
  \item Open the default address printed in the terminal and make it full screen to have a better view. 
\end{enumerate}

\section{Optional: production build}

\noindent
To run an optimised build (closer to a deployed Next.js app, useful for a stable demo):

\begin{verbatim}
npm run build
npm run start
\end{verbatim}

\section{Feature checklist (where to see each behaviour)}


\begin{center}
\small
\begin{tabular}{@{}P{3.6cm}P{9.8cm}@{}}
\hline
\textbf{Feature} & \textbf{What to do in the running prototype} \\
\hline
Pre-send permissions and screenshot warning & From the sender handset, open the photo picker, choose one or more images, then use \emph{Send with revocable permissions}. \\
\hline
Post-send permission changes / revocation & On the sender handset, open the action menu on a sent photo message and use the post-send options (e.g.\ revoke access, forwarding, download, time limit). Expand the permission summary row under the bubble if needed. \\
\hline
Recipient view and revoked placeholder & Compare the recipient handset before and after sender-side revocation. Use recipient menu actions such as \emph{Forward}, \emph{Download}, or \emph{Request access} when the design allows them. \\
\hline
View-once consumption & On the recipient handset, open a view-once photo via \emph{Open}, then close the viewer and observe the consumed state on the bubble. \\
\hline
Forwarding and access request & Forward from a recipient handset when permitted. After revocation, use \emph{Request access} and respond on the sender handset with \emph{Accept} or \emph{Reject}. Observe the forwarded-recipient panel for propagated state. \\
\hline
Co-subject flow & Use the co-subject handset panel. Read the notification card, open the photo menu, and use \emph{Request revocation} where available. Complete \emph{Accept}/\emph{Reject} on the sender handset. \\
\hline
\end{tabular}
\end{center}

\section{Optional: automated tests}

\noindent

\begin{verbatim}
npm run test:coverage
\end{verbatim}

\noindent
Further detail on testing is given in Section~\ref{sec:software-testing}.
